\section{Methodology: Design Science Research}

This research follows the \textit{Design Science Research} (DSR) methodology,
which provides a systematic approach for designing, developing, and evaluating
artifacts intended to address practical problems while contributing to
scientific knowledge. Unlike explanatory research approaches that primarily
aim to understand existing phenomena, DSR focuses on the creation and
evaluation of purposeful artifacts through an iterative design process.

DSR is appropriate for this thesis because the objective is to develop and
evaluate an AI-based regulatory document entity linking framework. Regulatory documents
contain numerous references to legal acts, directives, standards, and technical
specifications, where the same regulatory document entity may be expressed through
different identifiers, abbreviations, titles, or descriptive references.
Resolving these references to their corresponding canonical entities is a
fundamental requirement for constructing regulatory knowledge graphs and
enabling automated compliance workflows.

The research follows the Design Science Research Methodology (DSRM) proposed
by Peffers et al.~\cite{peffers2007design}, which consists of six iterative
activities:

\begin{enumerate}

    \item \textbf{Problem Identification and Motivation.}
    The research problem is identified through a review of regulatory
    compliance, entity linking, information retrieval, and natural language
    processing literature. This analysis reveals limitations in existing
    approaches for resolving regulatory references due to the domain-specific
    characteristics of legal documents, including heterogeneous reference
    formats, semantic ambiguity, cross-references between legal acts, and
    limited availability of labeled regulatory data.

    \item \textbf{Definition of Objectives.}
    Based on the identified problem, the objectives of the artifact are
    established. The primary objective is to design a framework capable of
    linking textual regulatory references to their corresponding canonical
    regulatory entities with high reliability. The framework should support
    semi-structured regulatory data, incorporate strategies for addressing
    limited annotations, and enable robust entity resolution across varying
    reference forms.

    \item \textbf{Design and Development.}
    The artifact developed in this research is an end-to-end regulatory document entity
    linking framework covering the stages from regulatory data preparation to
    model inference. The framework incorporates data preprocessing, data
    augmentation, semantic representation learning, candidate generation, entity
    resolution, and the creation of an inference pipeline for applying the
    developed models to unseen regulatory references. Transformer-based models
    are utilized to create contextual representations of regulatory references
    and candidate entities, while multi-stage retrieval approaches are
    investigated to balance retrieval efficiency and linking accuracy.
    Additionally, the relationship between Transformer architectures and
    associative memory models, particularly Modern Hopfield Networks, is examined
    as a theoretical perspective on the associative retrieval behaviour of these
    models.

    \item \textbf{Demonstration.}
    The developed framework is demonstrated on regulatory datasets containing
    references to legal entities. Demonstration scenarios focus on resolving
    different surface forms of regulatory references, including official
    identifiers, abbreviated names, and descriptive mentions, to their
    corresponding canonical regulatory entities.

    \item \textbf{Evaluation.}
    The developed artifact is evaluated using entity linking metrics and
    retrieval-based evaluation methods. The evaluation examines the framework's
    ability to correctly identify regulatory entities, resolve ambiguous
    references, and generalize across different regulatory reference patterns.
    The results are analyzed to determine the effectiveness of the proposed
    framework in supporting reliable automated compliance workflows.

    \item \textbf{Communication.}
    The artifact design, implementation process, evaluation results, and
    research findings are documented throughout this thesis. The conclusions
    summarize the contributions of the developed regulatory document entity linking
    framework, discuss limitations, and identify opportunities for future
    research in AI-supported regulatory knowledge management.

\end{enumerate}

The iterative nature of DSR allows insights obtained during implementation and
evaluation to guide refinements of the developed artifact. By combining the
design of a practical regulatory document entity linking framework with systematic
evaluation, this research contributes both an artifact for resolving regulatory
references and knowledge regarding the application of AI-based entity linking
methods in compliance-oriented domains.